\documentclass{article}
\usepackage{amsmath,amssymb,amsthm}
\newtheorem{theorem}{Theorem}
\begin{document}
\title{Paired exterior offsets: exact identities and a scoped obstruction}
\maketitle

\paragraph{Verification status.}
The scalar identities and the Farkas argument below are complete written
proofs.  The displayed leading coefficients in the sixteen-token
counterexample are checked by rational Laurent arithmetic, with guarded
precision and an explicitly nonsingular square system.  There is also a
standalone rational Farkas certificate at $\varepsilon=1/100$.
These calculations are not Lean-checked.  They do not prove either global
asymptotic target and change no definition of the original model.

\paragraph{An exterior pair identity.}
Let $D>x+\gamma$, with $\gamma>0$, and define
\[
 g_{D,\gamma,\kappa}(p,a)
 =\frac{a+\kappa}{D-p}-\frac{a-\kappa}{D+\gamma-p}.
\]
Then
\[
 g(x,\kappa)=g(x+\gamma,-\kappa)=\frac{2\kappa}{D-x}.
 \tag{1}
\]
Consequently two tokens $(x,\kappa),(x+\gamma,-\kappa)$ have equal
readouts for every exterior $D$, so their signed difference is zero under
either assignment to two input slots.  This is an infinite family of
exterior offsets for a fixed nonquadratic token pool.
The more general exact expansion is
\[
 g(p,a)=\frac{\gamma a}{(D-p)(D+\gamma-p)}
       +\kappa\left(\frac1{D-p}+\frac1{D+\gamma-p}\right).
 \tag{2}
\]
For $\kappa=\gamma$, its numerator is
$\gamma(a+2D+\gamma-2p)$.
All these identities follow by putting the fractions over a common
positive denominator.

\paragraph{The two decoder variants.}
Use the butterfly permutation $\pi_\theta$ and switch indices
$e=(k,s,g)$ defined in \texttt{butterfly-rectangle-decoder.tex}.
Put $H=2^d$ and
\[
 \gamma_k=\varepsilon^{2^k},\qquad
 p_y=\sum_{k=0}^{d-1}y_k\gamma_k,\qquad D_s=1-p_s.
\]
For an input slot $h$, the left mask is
$M_{k,g}(h)=1_{h\mathbin{\mathrm{div}}2^k=2g}$.
The Haar mask $U_{k,g}$ is $+1$ on that left child, $-1$ on its sibling,
and zero outside the parent block.  With $V=M$ or $V=U$, consider
\[
 F_{k,s,g}(a,\beta)=\beta_{k,g}
 +\sum_h V_{k,g}(h)
       g_{D_s,\gamma_k,\gamma_k}(p_{\pi_\theta(h)},a_h).
 \tag{3}
\]
The question is whether some real $a_h,\beta_{k,g}$ satisfy
$(1-2\theta_e)F_e>0$ for every switch.  Residues are fully target-adaptive.
Poles are the prescribed target permutation of the fixed hierarchical pool.

Both mask choices use two banks of $H$ separable affine-numerator
reciprocals before the existing axis and cube-admissibility transfers.
For $V=+1$, bank one has numerator $a+\kappa$ and denominator $D-p$;
bank two has numerator $-a+\kappa$ and denominator $D+\gamma-p$.
For $V=-1$, use numerators $a-\kappa,-a-\kappa$ and interchange those
denominators.  For $V=0$, use numerators $a-\kappa,-a+\kappa$ and give
both banks denominator $D+\gamma-p$.  Thus every numerator remains a
row term plus a query term.  No product of independently varying row and
query parameters is assumed to be affine.  Block biases have the already
checked $O(H)$ implementation.  Positivity on the entire local Boolean
cube still requires the established projective budget; a finite test at
one value of $\varepsilon$ does not establish that budget for arbitrarily
many rows.

\paragraph{Exact finite tests and the next gap.}
At $\varepsilon=1/100$, the left-mask variant is feasible for all $256$
upper-switch words when $H=8$, and passed the first $20$ reproducible
random sixteen-token words.  The Haar variant passed the same $20$
sixteen-token words.  Every reported feasible solve was followed by exact
rational substitution, not merely a solver status check.  Bottom-switch
bits only permute the two members of input pairs, preserving every higher
mask, and their singleton-query blocks admit arbitrary biases.
These finite successes do not establish a growing decoder.
In fact, the next theorem disproves universality of the left-mask variant
for all sufficiently small values of its hierarchy parameter.

\paragraph{Eliminating the block biases.}
For each block, let $s$ be a query with desired positive sign and $t$ one
with desired negative sign.  A block bias exists exactly when every
positive score exceeds every negative score.  Thus (3) is feasible
exactly when all reduced inequalities
\[
 L_{k,g,s,t}(a)=A_{k,g,s,t}\cdot a+B_{k,g,s,t}>0
 \tag{4}
\]
hold, where $A$ and $B$ are the differences of the two expressions in
(2), including the chosen masks.  Pure-sign blocks impose no reduced
inequalities.  A nonzero vector $\lambda\geq0$ with
$\sum\lambda A=0$ and $\sum\lambda B\leq0$ proves infeasibility by
summing (4).

\begin{theorem}[Failure of the prescribed left-mask hierarchy]
There is a fixed sixteen-token butterfly word $\theta$ and an
$\varepsilon_0>0$ such that, for every $0<\varepsilon<\varepsilon_0$,
no real residues or block biases satisfy (3) with $V=M$.
This statement allows arbitrary target-adaptive residues, but fixes the
pole prescription, the query times, and $\kappa_k=\gamma_k$ above.
\end{theorem}
\begin{proof}
The switch ordering is increasing $k$, then $s$, then $g$.
Take
\[
\theta=(0,0,0,0,0,0,0,0,\;
1,1,0,1,1,0,1,1,\;
0,0,1,1,0,0,0,0,\;
0,0,0,1,1,0,1,0).
\]
The following nine reduced inequalities suffice.  The table also gives
the first nonzero term of a positive rational-function multiplier.
\[
\begin{array}{c|c}
(k,g,s,t)&\lambda(\varepsilon)\text{ leading term}\\\hline
(1,1,1,0)&7\varepsilon^{26}\\
(2,0,2,1)&(7/3)\varepsilon^{21}\\
(2,0,3,1)&(7/3)\varepsilon^{20}\\
(3,0,0,6)&1\\
(3,0,1,6)&10\varepsilon^5\\
(3,0,2,4)&1\\
(3,0,2,6)&3\varepsilon\\
(3,0,5,3)&\varepsilon^3\\
(3,0,7,3)&\varepsilon
\end{array}
\tag{5}
\]
Here is an exact definition, independent of a numerical optimizer, of
these multipliers.  Set the sixth multiplier equal to one.  Solve the
eight homogeneous equations $\sum\lambda A_y=0$ for physical output
columns
\[
 y=0,1,2,5,7,11,12,14.
\]
These are all columns on which any of the nine selected masks is nonzero.
The eight by eight coefficient matrix is invertible over
$\mathbb Q(\varepsilon)$.  Valuation-pivoted elimination has successive
pivot leading terms
\[
2\varepsilon^3,\ -2\varepsilon^5,\ 3\varepsilon^8,
\ 2\varepsilon^9,\ -3\varepsilon^{12},\ -4\varepsilon^{15},
\ -5\varepsilon^{20},\ 12\varepsilon^{25}.
\tag{6}
\]
Thus the system defines rational functions uniquely.  Exact series
elimination gives (5) and
\[
 \sum\lambda B=-\frac83\varepsilon^{19}+O(\varepsilon^{20}).
 \tag{7}
\]
All leading coefficients in (5) are positive.  Equations (5)--(7)
therefore imply, for a common sufficiently small positive $\varepsilon$,
that every multiplier is positive and their affine weighted sum is
negative.  The residue coefficients cancel identically by their defining
eight equations; every other residue column is identically zero.
Summing the nine purported strict inequalities contradicts (7).
\end{proof}

\paragraph{Auditing the finite algebra.}
The checker \texttt{butterfly-paired-offset-laurent-check.py} constructs
(2)--(4) from the displayed word and computes (5)--(7) using rational
series.  It starts with $512$ coefficients, checks the precision before
every division, and retains at least $415$ coefficients.  Since every
pivot in (6) has a certified nonzero leading coefficient, the operations
are computations in the Laurent expansion of the uniquely defined
rational solution, not assertions based on approximate residuals.
The saved file \texttt{butterfly-paired-offset-h16-laurent.json} records
all these leading terms.  Independently,
\texttt{butterfly-paired-offset-h16-dual.json} gives a rational certificate
at $\varepsilon=1/100$; the main checker verifies it using only exact
rational arithmetic.  The analogous thirty-two-token certificate is
saved as well.

\paragraph{Scope.}
This counterexample does not cover Haar masks, different query times,
independently chosen per-head offsets, target-adaptive pole locations,
or a joint compiler whose query parameters are chosen after seeing the
whole sign table.  It is not an obstruction to the required reverse-axis
compiler or a lower bound on worst-case head complexity.  The distinct
paired-Haar construction and more general adaptive offset families remain
open.  The purpose of the counterexample is to prevent finite successful
samples from being mistaken for a universal hierarchy theorem.
\end{document}
